% Laget automatisk av elevark/lag_elevark.py fra opplegg/generell-relativitet.html. Endringer her blir overskrevet.
\documentclass[11pt]{article}
\usepackage{elevark}
\begin{document}
\elevarkhode{Fysikk 2 · Generell relativitetsteori}{Generell relativitetsteori – arbeidsark}{https://fysikklab.no/sim/ekvivalensprinsippet.html}{Ekvivalensprinsippet}

\section*{Oppgaver}

Oppgave 1–4 bruker Ekvivalensprinsippet\footnote{\url{fysikklab.no/sim/ekvivalensprinsippet.html}}, oppgave 5–6 bruker Lys som bøyes av tyngdekraften\footnote{\url{fysikklab.no/sim/lysavboyning.html}}, og oppgave 7–8 bruker GPS-klokkene\footnote{\url{fysikklab.no/sim/gps.html}}.

\begin{oppgaveliste}
\oppgave \fase{gjett}{Gjett}En ball kastes vannrett i en rakett som akselererer med 9,81~m/s² ute i rommet. Hvordan ser banen ut for en astronaut inne i raketten? Kjør «Akselererer i rommet» og sjekk.
\svar
\oppgave \fase{test}{Test}Sammenlign «I ro på jorda» og «Akselererer i rommet». Hvor langt kommer ballen før den treffer gulvet i hvert tilfelle? Hva er forskjellen sett utenfra?
\svar
\oppgave \fase{forklar}{Forklar}Sammenlign «Faller fritt» og «Svever i rommet». Forklar med ekvivalensprinsippet hvorfor astronautene på romstasjonen svever.
\svar
\oppgave \fase{forklar}{Forklar}Se på lyspulsen i «Akselererer i rommet», både innenfra og utenfra. Forklar hvorfor lyset ser ut til å bøyes innenfra. Hvorfor betyr det at lys også bøyes av tyngdekraften på jorda?
\svar
\oppgave \fase{test}{Regn}Regn ut avbøyningen $θ = 4γM/(Rc^{2})$ for lys som streifer solranden. Sola har massen $1{,}99 \cdot 10^{30}\,\mathrm{kg}$ og radius $6{,}96 \cdot 10^{8}\,\mathrm{m}$. Gi svaret i buesekunder, og sjekk med simuleringen.
\svar
\oppgave \fase{test}{Undersøk}Sammenlign avbøyningen ved overflaten av sola, en hvit dverg og en nøytronstjerne. Regn ut Schwarzschild-radien til sola. Hva ville skjedd om sola ble presset sammen til en kule med denne radien?
\svar
\oppgave \fase{test}{Undersøk}Velg GPS-satellitten. Hvor mye går klokka for fort eller for sakte per døgn på grunn av farten, og på grunn av tyngdefeltet? Hvilken effekt er størst?
\svar
\oppgave \fase{test}{Regn}Hvor stor posisjonsfeil gir 38,5~µs per døgn? Hvorfor må GPS-systemet ta hensyn til relativitetsteorien?
\svar
\oppgave \fase{utfordring}{Utfordring}En astronaut faller med føttene først mot et svart hull med massen 10 solmasser. Hvorfor blir hun strukket ut? Hvorfor er det verre ved et lite svart hull enn ved et stort?
\svar
\end{oppgaveliste}

\end{document}
